% 前置部分（封面、目录、摘要）
\ifShowCover % 如果在 metadata 里把封面开关打开
% 封面页：titlepage 环境会把该页作为独立封面，不受正文分页样式影响
\begin{titlepage} % 开始封面环境
    \pagenumbering{empty} % 封面页不显示页码
    \begin{center} % 封面整体内容居中
        \vspace*{2cm} % 顶部预留 2cm 空白
        {\heiti \zihao{1} \SchoolName}\\[1cm] % 学院名称（黑体一号），并在下方留 1cm
        {\heiti \zihao{1} \ReportType}\\[3cm] % 报告类型（黑体一号），并在下方留 3cm
        {\zihao{3} 课程名称: \underline{\quad \CourseName \quad}}\\[4cm] % 课程名称，下面划线用于视觉强调
        \zihao{3} % 后续封面表格统一使用三号字
        \begin{tabular}{ % 用表格做封面信息对齐
                >{\raggedright\arraybackslash}p{3cm} % 第1列：左对齐固定宽度 3cm
                >{\raggedright\arraybackslash}p{3.5cm} % 第2列：左对齐固定宽度 3.5cm
                >{\raggedright\arraybackslash}p{3cm} % 第3列：左对齐固定宽度 3cm
                >{\raggedright\arraybackslash}p{3.5cm} % 第4列：左对齐固定宽度 3.5cm
            }
            学生专业： & \StudentMajor & 学生班级： & \StudentClass \\ % 第一行：专业与班级
            学生学号： & \StudentID & 学生姓名： & \StudentName \\ % 第二行：学号与姓名
            指导教师： & \TeacherName & & \\ % 第三行：指导教师，后两格留空
            成绩评估： & \\[2cm] % 第四行：留出 2cm 空白给手写成绩
            评阅老师： & & 评阅时间： & \\ % 第五行：留空给评阅老师与时间
        \end{tabular}\\[4cm] % 结束封面信息表格，并在下方留 4cm
        {\kaishu \zihao{3} \today} % 封面底部日期（楷书三号）
    \end{center} % 结束居中环境
\end{titlepage} % 结束封面环境 % 就插入封面文件 cover.tex
\fi % 结束 ifShowCover 条件块

\pagenumbering{roman} % 把页码格式改为罗马数字（i, ii, iii...），常用于前置部分
\ifShowTOC % 如果目录开关打开
\tableofcontents % 自动生成目录
\newpage % 目录后强制换页
\fi % 结束 ifShowTOC 条件块

\pagenumbering{arabic} % 正文开始后把页码改回阿拉伯数字（1,2,3...）
\setcounter{page}{1} % 将页码计数器重置为 1
\maketitle % 输出题名区（标题/作者/单位），其样式在 report.sty 中重定义

\ifShowCnAbstract % 如果中文摘要开关打开
\renewcommand{\abstractname}{\zihao{5} \heiti \mdseries 摘\quad 要} % 自定义“摘要”标题样式
\begin{abstract} % 开始摘要环境
    \zihao{5} % 摘要正文使用五号字
    摘要内容。概括地陈述论文研究的目的、方法、结果、结论，要求200～300字。应排除本学科领域已成为常识的内容；不要把应在引言中出现的内容写入摘要，不引用参考文献；不要对论文内容作诠释和评论。不得简单重复题名中已有的信息。用第三人称，不使用“本文”、“作者”等作为主语。使用规范化的名词术语，新术语或尚无合适的汉文术语的，可用原文或译出后加括号注明。除了无法变通之外，一般不用数学公式和化学结构式，不出现插图、表格。缩略语、略称、代号，除了相邻专业的读者也能清楚理解的以外，在首次出现时必须加括号说明。结构严谨，表达简明，语义确切。 % 这里是中文摘要示例正文
    \par \noindent{{\heiti 关键词：}{关键词1；关键词2；关键词3；关键词4}} % 输出关键词并取消该行首行缩进
\end{abstract} % 结束摘要环境
\fi % 结束 ifShowCnAbstract 条件块

\ifShowEnAbstract % 如果英文摘要开关打开
\begin{center} % 以下内容居中排版
    \vspace{1em} % 在英文摘要标题前留一点竖直间距
    {\bf \zihao{3} \ReportNameEn \par\vspace*{10pt}} % 英文题目（加粗+三号字）并在后面留间距
    {\zihao{-4} \StudentNameEn}\\ % 英文作者名，\\ 表示换行
    {\zihao{5} \StudentAffiliationEn} % 英文单位
\end{center} % 结束居中环境
\renewcommand{\abstractname}{\zihao{5} \bf Abstract} % 把摘要标题改成英文 Abstract
\begin{abstract} % 开始英文摘要环境
    \zihao{5} % 英文摘要正文用五号字
    Purpose purpose purpose purpose purpose purpose purpose purpose purpose purpose purpose purpose purpose purpose purpose purpose purpose purpose purpose purpose purpose purpose purpose purpose purpose purpose purpose purpose purpose purpose purpose purpose. Method method method method method method method method method method method method method method method method method method. Result result result result result result result result result result result result result result result result result result result result result result result result result result result. Conclusion conclusion conclusion conclusion conclusion conclusion conclusion conclusion conclusion conclusion conclusion conclusion conclusion conclusion conclusion conclusion conclusion conclusion conclusion conclusion. % 这里是英文摘要占位示例
    \par \noindent{\textbf{Keywords:}{keyword1; keyword2; keyword3; keyword4}} % 英文关键词行
\end{abstract} % 结束英文摘要环境
\fi % 结束 ifShowEnAbstract 条件块